% -----------------------------------------------------------------------------
% R2: rotational relative-equilibrium residual class
% -----------------------------------------------------------------------------
% This section closes the rotational relative-equilibrium class without using a
% whole-space Coriolis--Leray Liouville theorem.  The only estimates introduced
% here are local.  The Coriolis contribution is recorded as a localized annular
% flux, and every non-small annular flux is routed into the same retained-active
% local concentration alternatives used later in the terminal analysis.
% -----------------------------------------------------------------------------

\section{The rotational relative-equilibrium residual class}
\label{sec:R2-reduction}

We next isolate the second refined residual class.  This class consists of
centered Type~I ancient limits whose time dependence is a rigid rotation of a
fixed spatial profile.  Passing to the co-rotating frame converts the centered
Navier--Stokes equation into a stationary Leray profile equation with a
Coriolis transport term.  Earlier drafts closed this class by postulating a
bounded whole-space Coriolis--Leray Liouville theorem.  We do not use that
global statement here.  Instead, the Coriolis term is localized as an annular
flux and is then treated by the local concentration analysis proved in
\Cref{sec:terminal-concentration,sec:separated-profiles-atomic}.

Throughout this section the centered equation is
\begin{equation}\label{eq:centered-R2-safe}
   \partial_\tau V-\Delta V+\frac12 y\cdot\nabla V+\frac12 V
   +(V\cdot\nabla)V+\nabla P=0,
   \qquad \nabla\cdot V=0 .
\end{equation}
For \(\Omega\in\mathfrak{so}(3)\), set
\begin{equation}\label{eq:rotation-group-R2-safe}
   Q(\tau):=\exp(\tau\Omega)\in SO(3).
\end{equation}

\begin{definition}[Rotational relative-equilibrium profile]
\label{def:R2-relative-profile-safe}
A bounded smooth ancient solution \((V,P)\) of \eqref{eq:centered-R2-safe} is a
rotational relative equilibrium if there exist
\(\Omega\in\mathfrak{so}(3)\), a bounded smooth divergence-free vector field
\(W\), and a pressure profile \(\Pi\), such that
\begin{equation}\label{eq:R2-relative-form-safe}
   V(y,\tau)=Q(\tau)W(Q(\tau)^Ty),
   \qquad
   P(y,\tau)=\Pi(Q(\tau)^Ty)
\end{equation}
up to the usual time-dependent pressure gauge.
The class \(\Crot\) consists of such residual Seregin profiles, after the
axisymmetric bounded-circulation class and the previously closed Paper I--III
alternatives have been removed.
\end{definition}

\begin{lemma}[Co-rotating stationary equation]
\label{lem:co-rotating-equation-safe}
Let \((V,P)\) be a rotational relative equilibrium.  Then the co-rotating
profile \((W,\Pi)\) satisfies
\begin{equation}\label{eq:R2-coriolis-profile-safe}
   -\Delta W+\frac12 y\cdot\nabla W+\frac12 W
   +(W\cdot\nabla)W
   +\Omega W-(\Omega y)\cdot\nabla W
   +\nabla\Pi=0,
   \qquad \nabla\cdot W=0 .
\end{equation}
Equivalently,
\begin{equation}\label{eq:R2-coriolis-profile-project-safe}
   \mathbb P\Big[
   -\Delta W+\frac12 y\cdot\nabla W+\frac12 W
   +(W\cdot\nabla)W
   +\Omega W-(\Omega y)\cdot\nabla W
   \Big]=0,
\end{equation}
where \(\mathbb P\) is the whole-space Leray projector.
\end{lemma}

\begin{proof}
Let \(z=Q(\tau)^Ty\).  Since \(Q'(\tau)=\Omega Q(\tau)\) and
\(Q(\tau)^T\Omega Q(\tau)=\Omega\), one has
\[
   \partial_\tau z=-\Omega z.
\]
Therefore
\[
   \partial_\tau V(y,\tau)
   =Q(\tau)\Big(\Omega W(z)-(\Omega z)\cdot\nabla W(z)\Big).
\]
The Laplacian, divergence, convection term, and pressure gradient are invariant
under the orthogonal change of variables.  Substitution into
\eqref{eq:centered-R2-safe}, followed by multiplication by \(Q(\tau)^T\), gives
\eqref{eq:R2-coriolis-profile-safe}.  Applying the Leray projector gives
\eqref{eq:R2-coriolis-profile-project-safe}.
\end{proof}

\subsection{The Bernoulli identity and the local Coriolis flux}

Write
\begin{equation}\label{eq:R2-F-def}
   F(y):=\frac12y-\Omega y,
   \qquad
   A(y):=W(y)+F(y),
   \qquad
   \omega:=\nabla\times W .
\end{equation}
Since \(\Omega\) is skew-symmetric, there is a unique vector
\(\varpi\in\mathbb R^3\) such that
\begin{equation}\label{eq:R2-Omega-vector}
   \Omega y=\varpi\times y .
\end{equation}
Then
\begin{equation}\label{eq:R2-curlF}
   \nabla\times F=-2\varpi .
\end{equation}

\begin{lemma}[Coriolis Bernoulli identity]
\label{lem:R2-coriolis-bernoulli}
Let \((W,\Pi)\) be a smooth solution of
\eqref{eq:R2-coriolis-profile-safe}.  Define
\begin{equation}\label{eq:R2-Bernoulli-def}
   B:=\frac12|W|^2+\Pi+F\cdot W .
\end{equation}
Then
\begin{equation}\label{eq:R2-gradient-B}
   \nabla B+\nabla\times\omega-A\times\omega=0,
\end{equation}
and consequently
\begin{equation}\label{eq:R2-Bernoulli-PDE}
   -\Delta B+A\cdot\nabla B
   =
   -|\omega|^2-\omega\cdot(\nabla\times F)
   =
   -|\omega|^2+2\varpi\cdot\omega .
\end{equation}
\end{lemma}

\begin{proof}
The vector identity
\[
   (W\cdot\nabla)W=\nabla\left(\frac12|W|^2\right)-W\times\omega
\]
and the identity
\[
   (F\cdot\nabla)W+(\nabla F)^TW
   =\nabla(F\cdot W)-F\times\omega
\]
apply because
\((\nabla F)^TW=\frac12W+\Omega W\).  Thus
\eqref{eq:R2-coriolis-profile-safe} can be rewritten as
\[
   -\Delta W+
   \nabla\left(\frac12|W|^2+\Pi+F\cdot W\right)
   -(W+F)\times\omega=0 .
\]
Since \(\nabla\cdot W=0\), one has
\(-\Delta W=\nabla\times\omega\), which gives
\eqref{eq:R2-gradient-B}.

Taking divergence in \eqref{eq:R2-gradient-B} gives
\[
   \Delta B=\nabla\cdot(A\times\omega),
\]
because the divergence of a curl vanishes.  Using
\[
   \nabla\cdot(A\times\omega)
   =\omega\cdot(\nabla\times A)-A\cdot(\nabla\times\omega),
   \qquad
   \nabla\times A=\omega+\nabla\times F,
\]
we obtain
\[
   \Delta B
   =|\omega|^2+\omega\cdot(\nabla\times F)
      -A\cdot(\nabla\times\omega).
\]
On the other hand, by \eqref{eq:R2-gradient-B},
\[
   A\cdot\nabla B
   =A\cdot(A\times\omega)-A\cdot(\nabla\times\omega)
   =-A\cdot(\nabla\times\omega).
\]
Combining the last two identities gives \eqref{eq:R2-Bernoulli-PDE}.  Finally,
\eqref{eq:R2-curlF} follows from \eqref{eq:R2-Omega-vector}.
\end{proof}

\begin{lemma}[Localized Coriolis flux identity]
\label{lem:R2-coriolis-flux-localization}
Let \((W,\Pi)\) be as in \Cref{lem:R2-coriolis-bernoulli}, and let
\(\chi\in C_c^\infty(\mathbb R^3)\).  Then
\begin{equation}\label{eq:R2-local-coriolis-source}
   \int_{\mathbb R^3}\chi
   \left(
      -\Delta B+A\cdot\nabla B+|\omega|^2
   \right)\,dy
   =
   -2\int_{\mathbb R^3}(W\times\varpi)\cdot\nabla\chi\,dy .
\end{equation}
In particular the Coriolis contribution is supported only where the cutoff
varies.
\end{lemma}

\begin{proof}
By \eqref{eq:R2-Bernoulli-PDE},
\[
   -\Delta B+A\cdot\nabla B+|\omega|^2=2\varpi\cdot\omega .
\]
Because \(\varpi\) is constant and \(\omega=\nabla\times W\),
\[
   \varpi\cdot\omega
   =
   \nabla\cdot(W\times\varpi).
\]
Multiplying by \(\chi\), integrating over \(\mathbb R^3\), and integrating by
parts gives
\[
   2\int\chi\,\varpi\cdot\omega
   =
   2\int\chi\,\nabla\cdot(W\times\varpi)
   =
   -2\int(W\times\varpi)\cdot\nabla\chi ,
\]
which proves \eqref{eq:R2-local-coriolis-source}.
\end{proof}

\begin{lemma}[Non-small Coriolis flux creates an active local concentration]
\label{lem:R2-annular-flux-active-concentration}
Let \(\chi\in C_c^\infty(\mathbb R^3)\), let
\(A_\chi:=\operatorname{supp}\nabla\chi\), and put
\[
   C_\chi:=\int_{A_\chi}|\nabla\chi|^{3/2}\,dy .
\]
Fix \(T\in(0,1]\), and let
\[
   \mathcal T_{\chi,T}
   :=
   \{(y,\tau): -T<\tau<0,\ Q(\tau)^Ty\in A_\chi\}
\]
be the swept annular tube in the original centered variables.  Cover
\(\mathcal T_{\chi,T}\) by finitely many unit terminal cylinders
\[
   \mathcal T_{\chi,T}\subset \bigcup_{m=1}^{N_{\chi,T}}Q_1(z_m).
\]
Let \(\eps_{\rm sd}>0\) be the local activity threshold used in
\Cref{def:terminal-measures}.  If \(\varpi\ne0\) and
\begin{equation}\label{eq:R2-flux-active-threshold}
   \left|
      2\int_{\mathbb R^3}(W\times\varpi)\cdot\nabla\chi\,dy
   \right|
   \ge
   2|\varpi|\,C_\chi^{2/3}
      \left(\frac{N_{\chi,T}\eps_{\rm sd}}{T}\right)^{1/3},
\end{equation}
then at least one cylinder \(Q_1(z_m)\) is active for the original
centered solution \(V\).
\end{lemma}

\begin{proof}
Write
\[
   \delta_\chi:=
   \left|
      2\int_{\mathbb R^3}(W\times\varpi)\cdot\nabla\chi\,dy
   \right| .
\]
Holder's inequality gives
\[
\begin{aligned}
   \delta_\chi
   &\le
   2|\varpi|\int_{A_\chi}|W|\,|\nabla\chi|\,dy  \\
   &\le
   2|\varpi|
   \left(\int_{A_\chi}|W|^3\,dy\right)^{1/3}
   C_\chi^{2/3}.
\end{aligned}
\]
Thus \eqref{eq:R2-flux-active-threshold} implies
\[
   \int_{A_\chi}|W|^3\,dy\ge \frac{N_{\chi,T}\eps_{\rm sd}}{T}.
\]
Since \(V(y,\tau)=Q(\tau)W(Q(\tau)^Ty)\) and \(Q(\tau)\) is an isometry,
\[
\begin{aligned}
   \iint_{\mathcal T_{\chi,T}} |V(y,\tau)|^3\,dy\,d\tau
   &=
   \int_{-T}^{0}\int_{Q(\tau)A_\chi}
      |W(Q(\tau)^Ty)|^3\,dy\,d\tau \\
   &=
   T\int_{A_\chi}|W(z)|^3\,dz
   \ge N_{\chi,T}\eps_{\rm sd}.
\end{aligned}
\]
Because the tube is covered by \(N_{\chi,T}\) unit terminal cylinders, one of
them satisfies
\[
   \iint_{Q_1(z_m)} |V|^3\,dy\,d\tau\ge \eps_{\rm sd}.
\]
The pressure contribution to the local CKN measure is nonnegative after the
allowed pressure gauge.  Hence
\[
   \mu(Q_1(z_m))\ge \eps_{\rm sd},
\]
so the corresponding recentering sequence is active in the sense of
\Cref{def:terminal-measures}.
\end{proof}

\begin{lemma}[Local Coriolis-flux routing dichotomy]
\label{lem:R2-flux-routing-dichotomy}
Fix \(\chi\in C_c^\infty(\mathbb R^3)\) and \(T\in(0,1]\).  Along every
terminal extraction of a rotational relative equilibrium, exactly one of the
following local alternatives occurs:
\begin{enumerate}[label=(\roman*)]
   \item \(\varpi\ne0\), the flux across
   \(A_\chi=\operatorname{supp}\nabla\chi\) satisfies the threshold
   \eqref{eq:R2-flux-active-threshold}, and the swept annular tube contains an
   active unit-scale concentration;
   \item either \(\varpi=0\), or the threshold
   \eqref{eq:R2-flux-active-threshold} fails; then the flux is a bounded
   compactly supported functional of \(W\) on \(A_\chi\), passes to terminal
   limits by local \(L^3\) convergence, and creates no terminal state other than
   the standard single-profile, separated-profile, exterior, radiative,
   scattering, or rough
   alternatives of \Cref{thm:terminal-exhaustion-main}.
\end{enumerate}
\end{lemma}

\begin{proof}
Alternative \emph{(i)} is precisely
\Cref{lem:R2-annular-flux-active-concentration}.  If \(\varpi=0\), then the flux
vanishes identically.  If \(\varpi\ne0\) and
\eqref{eq:R2-flux-active-threshold} fails, the flux is bounded by the right hand
side of that threshold.  In both cases the functional
\[
   W\longmapsto -2\int_{\mathbb R^3}(W\times\varpi)\cdot\nabla\chi\,dy
\]
depends only on \(W\) on the compact annular set \(A_\chi\).  Under the local
Seregin compactness supplied by Paper I, velocities converge strongly in
\(L^3_{\rm loc}\); since \(\nabla\chi\in L^{3/2}\) and \(\varpi\) is fixed, the
functional is continuous along every terminal subsequence:
\[
   \int (W_n\times\varpi)\cdot\nabla\chi
   \longrightarrow
   \int (W_\infty\times\varpi)\cdot\nabla\chi .
\]
Thus a subthreshold flux is retained only as a local term in the limiting
localized Bernoulli identity \eqref{eq:R2-local-coriolis-source}.  It is not a
new carrier of local velocity mass.  Every actual carrier of compact velocity
\(L^3\)-mass is still one of the terminal local concentration alternatives listed in
\Cref{thm:terminal-exhaustion-main}.
\end{proof}

\begin{remark}[Why this is not a global Liouville argument]
\label{rem:R2-no-global-liouville}
For \(\Omega=0\), the Coriolis flux in
\Cref{lem:R2-coriolis-flux-localization} vanishes.  For \(\Omega\ne0\), the
same lemma does not give a sign-definite maximum principle for \(B\).  It gives
only the local dichotomy in \Cref{lem:R2-flux-routing-dichotomy}: either a
chosen annulus already contains an active recentered local concentration, or the
annular flux is a subthreshold compact functional which passes to the local
terminal limit.  No tail decay, whole-space \(L^p\) bound, Gaussian energy
estimate, or global pressure normalization is used.
\end{remark}

\subsection{Local closure of the rotational relative-equilibrium class}

\begin{theorem}[Stratified exclusion of rotational relative equilibria]
\label{thm:R2-exclusion-safe}
The ordered rotational relative-equilibrium residual class is empty:
\begin{equation}\label{eq:R2-empty-safe}
   \Crot(\mathcal S;I,J)=\varnothing .
\end{equation}
\end{theorem}

\begin{proof}
Suppose, for contradiction, that
\(V\in\Crot(\mathcal S;I,J)\).  By
\Cref{def:R2-relative-profile-safe} and
\Cref{lem:co-rotating-equation-safe}, there are
\(\Omega\in\mathfrak{so}(3)\) and a bounded smooth divergence-free profile
\((W,\Pi)\) satisfying the Coriolis--Leray equation
\eqref{eq:R2-coriolis-profile-safe}.  The centered representative \(V\) is a
normalized Seregin ancient limit, so by
\Cref{hyp:base-seregin-package} it carries a fixed positive compact
velocity \(L^3\)-mass.

Choose a countable family of smooth cutoffs which is cofinal among compact
annuli in the terminal variables.  For each cutoff and each rational
\(T\in(0,1]\), apply the local dichotomy
\Cref{lem:R2-flux-routing-dichotomy}.  A superthreshold Coriolis flux produces
an active unit-scale concentration and is therefore already part of the active
local concentration decomposition.  A subthreshold Coriolis flux is a compactly
supported functional
of \(W\), continuous under the local convergence supplied by Paper I, and hence
remains inside the localized limiting equation; it carries no independent CKN
mass and creates no additional terminal stratum.

We now apply the terminal local concentration decomposition to the inherited compact
density.  The rough, inactive, exterior, scattering, and radiative alternatives
cannot be the sole carriers of this density by
\Cref{thm:terminal-nonactive-discharge}.  By the preceding paragraph, the
Coriolis flux has already been routed either into an active local concentration or into a
compact local term in the selected terminal equation.  Thus, after local active
extraction, the retained density is carried by either a finite separated family
of retained active profiles or a singleton retained active terminal profile.

The finite separated-family alternative is impossible by
\Cref{thm:no-finite-separated-profile-family}.  In the singleton case,
\Cref{prop:no-multi-implies-atomic} makes the terminal profile atomic, and
\Cref{lem:no-atomic-active} excludes retained active atomic profiles.  Both
possible carriers of the inherited compact density are impossible.  This
contradicts the normalization of \(V\), and therefore
\(\Crot(\mathcal S;I,J)=\varnothing\).
\end{proof}

\begin{proposition}[Already-closed integrable subcase]
\label{prop:R2-integrable-subcase}
Let \(V\) be a rotational relative equilibrium with co-rotating profile \(W\).
If
\begin{equation}\label{eq:R2-L3-subcase}
   W\in L^3(\mathbb R^3),
\end{equation}
then \(V\equiv0\).  More generally, the same conclusion holds under any
condition which places the physical representative in
\(L^\infty_tL^3_x\) on its ancient time interval.
\end{proposition}

\begin{proof}
Rotations and the Navier--Stokes critical scaling preserve the \(L^3\) norm.
Thus \(W\in L^3\) implies
\[
   \sup_{\tau\in\mathbb R}\|V(\cdot,\tau)\|_{L^3(\mathbb R^3)}<\infty .
\]
Equivalently, the physical representative lies in the critical class
\(L^\infty_tL^3_x\).  The endpoint \(L^3\) class is one of the already closed
alternatives in Paper I, by the Escauriaza--Seregin--\v Sver\'ak and stationary
Ne\v cas--R\r u\v zi\v cka--\v Sver\'ak inputs cited there.  Therefore
\(V\equiv0\).  A zero profile has zero local velocity mass, so no such profile
can belong to the residual
nonvanishing class.
\end{proof}
